\scanpage{26}
\mspage{13}
considérer que \deleted{\unsure{le|que}} $\tau$ définit \inserted{(grâce à l'op. de $\pi$)} un foncteur
\diagram{\begin{tikzcd}[ampersand replacement=\&,row sep=large,column sep=huge]
\operatorname{Rev}(\Pi) \arrow[r,"F"] \& (\operatorname{Ens})\\
(\pi\text{-}\operatorname{Ens}) \arrow[u,"\simeq"',"\begin{gathered}\deleted{\text{via }\widetilde\Pi}\\E\longmapsto\widetilde\Pi^{\circ}\times_\pi E\end{gathered}"] \arrow[ur,bend right=25,"E\longmapsto\tau^{\circ}\wedge_\pi E"']
\end{tikzcd}}
[\uline{NB} l'indice $\circ$ indique qu'on a renversé les op. de $\pi$, pour faire opérer $\pi$ à dr)
foncteur commutant aux $\varinjlim$ quelconques. Ceci dit, on veut que $\unsure{\Gamma|F}$ (qui opère sur $\operatorname{Rev}(\Pi)$ via ses opérations sur $\Pi$) soit compatible avec le foncteur $F$, i.e. qu'on ait des isom
\[
F({}^{\gamma}\Pi')\simeq F(\Pi')
\]
$(\Pi'\in\operatorname{Ob}\operatorname{Rev}\Pi)$ fonctoriels en $\Pi'$, et avec transitivité pour $\gamma$ variable. On trouve que celà \unsure{s'exprime}
\mspage{25}
